\documentclass[a4paper, 11pt]{ltjsarticle}
\usepackage[margin=15mm]{geometry}
\usepackage{amsmath}
\usepackage{multicol}

\usepackage{bm}
\usepackage{array}
\usepackage{calc}
\usepackage[table]{xcolor} 
\usepackage{makecell}      

\setlength{\parindent}{0pt}
\setlength{\columnsep}{25mm}
\newlength{\probheight}

% --- 枠付き筆算用マクロ(2x1用) ---
\newcommand{\hissangrid}[5]{
    \begin{minipage}[t][\probheight][t]{\linewidth}
        \vspace*{2ex}
        {\Large #1)} \par
        \vspace*{1em}
        \begin{center}
        \renewcommand{\arraystretch}{2.0}
        \arrayrulecolor{black!30}
        \begin{tabular}{|w{c}{2.2em}|w{c}{2.2em}|w{c}{2.2em}|w{c}{2.2em}|}
            \hline
            & & #2 \\ \hline 
            & \color{black}#5 & & #3 \\ 
            \noalign{\global\arrayrulewidth=1.5pt}\arrayrulecolor{black}\hline
            \noalign{\global\arrayrulewidth=0.4pt}\arrayrulecolor{black!30}
            & #4 \\ \hline
        \end{tabular}
        \arrayrulecolor{black}
        \end{center}
    \end{minipage}
}

\pagestyle{empty}

\begin{document}

% \hskip の代わりに、より安全な \hspace{1zw} を使用します
% \hspace{1em} {\Large \textbf{筆算トレーニング（2桁×1桁）}}
\hfill

% \rightline{　年\quad 　組\quad 　番 \quad 氏名\underline{\hspace{25ex}}}
次の計算をしなさい。
% 名前欄と問題の間に 5ex の間隔を空ける
\vspace*{5ex}

\begin{multicols*}{2}
\setlength{\probheight}{(\textheight - 60mm) / 3}
\hissangrid{1}{7 & 1}{4}{ & & }{×} 

\hissangrid{2}{6 & 9}{9}{ & & }{×} 

\hissangrid{3}{9 & 4}{6}{ & & }{×} 

\columnbreak
\hissangrid{4}{4 & 7}{3}{ & & }{×} 

\hissangrid{5}{3 & 5}{6}{ & & }{×} 

\hissangrid{6}{3 & 4}{1}{ & & }{×} 

\end{multicols*}
\newpage
\setlength{\probheight}{(\textheight - 20mm) / 3}
\begin{multicols*}{2}
\hissangrid{7}{4 & 5}{8}{ & & }{×} 

\hissangrid{8}{5 & 5}{4}{ & & }{×} 

\hissangrid{9}{3 & 2}{5}{ & & }{×} 

\columnbreak
\hissangrid{10}{1 & 5}{3}{ & & }{×} 

\hissangrid{11}{1 & 3}{2}{ & & }{×} 

\hissangrid{12}{9 & 3}{6}{ & & }{×} 


\end{multicols*}

\end{document}